% !TEX root = 00_arbeit.tex

%---------------------------------------------------------------------------------
%% Einstellungen die am Anfang des Dokumentes erfolgen

%-------------------------------------------
%% Zusätzliche Pakete werden geladen
%\documentclass[a4paper,12pt,titlepage]{article}  
\documentclass[a4paper, 12pt, titlepage]{scrartcl} % Schrift- und Papiergröße werden eingestellt  
\usepackage[top=2.5cm, bottom=2cm, left=2.5cm, right=2.5cm]{geometry} % Korrekturrand wird eingestellt
\usepackage[ngerman]{babel}             % Dokumentsprache wird auf deutsch gesetzt
\usepackage[table]{xcolor}              % Ermöglicht die Änderung der Schriftfarbe
\usepackage{graphicx}                   % Ermöglicht das Einbinden von Bildern
\usepackage{fancyhdr}                   % Ermöglicht die freie Bearbeitung der Fuß und Kopfzeilen
\usepackage{tocloft}                    % Ermöglicht das anpassen des Inhaltsverzeichnisses
%\usepackage{wrapfig}                    % Ermöglicht textumschlossene Abbildungen
\usepackage{array}                      % Ermöglicht das Anpassen von Tabellen
\usepackage{amsmath}                    % Erweitert die Mathematikumgebung
\usepackage{ifthen}                     % Ermöglicht wenn dann Abfragen
\usepackage{verbatim}                   % Ermöglicht das Auskommentieren mehrerer Zeilen an Code
%\usepackage{caption,subcaption}
%\usepackage{subfig}                     % Ermöglicht mehrere Bilder in  einer Abbildung
%\usepackage{fp}                         % Ermöglicht Kalkulationen innerhalb von LaTeX


%-------------------------------------------
%% Stellt Zeilenabstand auf Word vergleichbaren Abstand von 1,5 
%\usepackage[onehalfspacing]{setspace}   % Zeilenabstand wird auf 1,5 gesetzt
\usepackage{setspace}
%\FPeval\zeilenabstand{175/122}            % Berechnet den Wert innerhalb der Klammern und speichert ihn in der Variable \zeilenabstand
\setstretch{1.428}



%-------------------------------------------
%% Fonteinstellungen abhängig vom Compiler
%% Quellen: https://en.wikibooks.org/wiki/LaTeX/Fonts
%%          https://www.sharelatex.com/learn/XeLaTeX
%%          https://tex.stackexchange.com/questions/12565/load-fonts-that-are-in-a-fonts-directory
\usepackage{ifxetex}                    % Ermöglicht das Prüfen ob XeLaTeX benutz wird
    \ifxetex
        \usepackage{fontspec}
        \defaultfontfeatures{Ligatures=TeX} % To support LaTeX quoting style
        \setmainfont[%                      % Times New Roman (aus Windows) wird als Standardschrift gesetzt
            Path = Fonts/,%                 % Dateien Pfad wird auf Fonts/ gesetzt
            BoldFont=timesbd.ttf,%
            ItalicFont=timesi.ttf,%
            BoldItalicFont=timesbi.ttf,%
            ]{times.ttf}
    \else
        \usepackage[T1]{fontenc}            % Liefert die meisten Text-Zeichen für westeuropäische Sprachen
        \usepackage[utf8]{inputenc}         % Ermöglicht die direkte Eingabe der Umlaute
        \usepackage{mathptmx}               % Times wird als Standardschrift gesetzt
    \fi

%-------------------------------------------
%% Pakete für das Erstellen von Refferenzen
\usepackage{varioref}
\usepackage{hyperref} % Erweitert die Referenzmöglichkeiten
\usepackage{cleveref}

%-------------------------------------------
%% Pakete für das Erstellen von Tabellen
%\usepackage{tabu}
%\usepackage{longtable}
\usepackage{tabularx}
\usepackage{booktabs}
\newcolumntype{Z}{>{\raggedright\let\newline\\\arraybackslash}X}    % Ermöglicht linksbündige Spalten mit Zeilenumbruch in tabularx durch verwenden von Z 
\definecolor{tableShade}{gray}{0.95}     % Einfärben der Zeilen in Grau



%-------------------------------------------
%% Pakete für das Erstellen von Grafiken
\usepackage{tikz-cd}                    % Ermöglicht das Erstellen von Grafiken
\usepackage{circuitikz}                 % Ermöglicht das Zeichnen elektrischer Schaltungen und logischer Blöcke 
%\usepackage{relsize}                   % Ermöglicht die Fontgröße zu skalieren
\usetikzlibrary{arrows,calc,intersections,shapes}

%-------------------------------------------
%% Pakete für die richtige Darstellung von Fußnoten
\usepackage[bottom]{footmisc}           % Setzt die Fußnoten ans Ende der Seite
\usepackage{caption}
%\usepackage{bigfoot}
%\usepackage{savefnmark}
%\footnotesep\baselineskip
%\deffootnote{0.6cm}{1em}{\makebox[0.6cm][l]{\thefootnotemark}}
\setkomafont{footnote}{\fontsize{10pt}{10.5pt}\selectfont}  % Stellt die Fußnotengröße auf 10 ein
\setlength{\footnotesep}{13.5pt}                              % Sellt den Fußnotenzeilenabstand auf einfachen Zeilenabstand ein

%-------------------------------------------
%% Pakete für das Erstellen des Literaturverzeichnisses
\usepackage[babel,german=quotes]{csquotes}  % Verwendet die deutsche Art der Anführungszeichen
\usepackage[backend=biber, style=authoryear, citestyle=authoryear, maxitems=1]{biblatex}

\DefineBibliographyStrings{german}{%    % Verändert beim zitieren "u.a." zu "et al." und "und" zu "and"
  andothers = {et al.},
  %and       = {and}
}
\addbibresource{11_literatur.bib}       % Fügt die Literaturverzeichnisdatei ein
\setlength{\bibitemsep}{10pt}           % Stellt den Zeilenabstand zwischen den Literaturverzeichnis Einträgen ein
%-------------
%% Erstellt ein Makro welches prüft ob die eingegebene Variable leer ist
\def\ifempty#1{% 
\def\temp{#1}%
\ifx\temp\empty%
}

%-------------
%% Erstellt mit \citef ein Zitat in der Fußzeile in der Form "Vgl. Kramer (2009), S. 2."
\newcommand{\citef}[2][]{%              
    \ifempty{#1}%
        \footnote{\,Vgl.~\citeauthor{#2}~(\citeyear{#2}).}%
    \else%
        \footnote{\,Vgl.~\citeauthor{#2}~(\citeyear{#2}),~S.~#1.}%
    \fi%
}
%-------------
%% Statt \textcite kann nun auch \citet verwendet werden
\newcommand{\citet}[2][]{%              
    \ifempty{#1}%
        \citeauthor{#2}~(\citeyear{#2})%
    \else%
        \citeauthor{#2}~(\citeyear{#2}),~S.~#1%
    \fi%
}
%-------------
%% Erstellt mit \citemp ein Platzhalterzitat
\newcommand\citemp{%              
\def\footnotelayout{\color{red}}%
        \footnote{\farbig{Vgl.~Platzhalter~(0000),~S.~1.}}%
}



%-------------------------------------------
%% Erstellt eine TODO-Box mit \todo{}

\newcommand{\farbig}[2][red]{\textcolor{#1}{#2}}      % Ermöglicht text mit \farbig{text} rot einzufärben

\newcommand{\todo}[1]{
\begin{figure}[!htb]
    \centering
    \tikz \node[draw,text width=.8\textwidth,align=left]{\textbf{TODO:}\\\farbig{#1.}};
    %\fbox{\parbox{.8\textwidth}{\textbf{TODO:}\\\farbig{#1.}}}
\end{figure}
}


%-------------------------------------------
\usepackage{lipsum}                     % Erstellt blindtext

%-------------------------------------------
%% Deckblatt
\usepackage{pdfpages}                   % Ermöglicht das Einbinden von PDF Dokumenten (wird verwendet für die Einbindung der Aufgabenstellung)
\usepackage[showboxes,absolute]{textpos}   % Ermöglicht das Setzen von Inhalt an einem beliebigen Ort
\usepackage{tikz}

\newcommand\Mygrid{%                        % Erzeugt ein Grid für textpos mit \Mygrid        
\tikz[
  remember picture,
  overlay,
  yscale=-1,
  xstep=\TPHorizModule,ystep=\TPVertModule,
  yshift=0pt,xshift=0pt]
  \draw (current page.north west) grid (current page.south east);}

%  yshift=13pt,xshift=4pt]
%-------------------------------------------

%-------------------------------------------
%% Erstellt inline Tikz Bilder
\newcommand{\inlineneuron}[2][1]{
    \tikz[baseline=-0.5ex,scale=#2, every node/.style={transform shape}]{
        \def\siz{.8}                % Größe der Neuronen
        \tikzstyle{neuron}=[circle,draw=black,minimum size=\siz cm,inner sep=2pt]
        \tikzset{
            ident/.pic={
            \draw[semithick] (-\siz/4,-\siz/4) -- (\siz/4,\siz/4);
        }}%
        \tikzset{
            ident8/.pic={
            \draw[semithick] (-\siz/8,-\siz/8) -- (\siz/8,\siz/8);
        }}%
%------------- Eingabeneuron -------------%
            \ifthenelse{\equal{#1}{i}}{%
                \node[neuron] (Ir) at (0,0) {};
                \pic at (Ir) {ident};
            }{}%
%--------------- Perceptron --------------%
            \ifthenelse{\equal{#1}{per}}{%
                \node[neuron,circle split,inner sep=2pt] (Or) at (0,0)
                    {\includegraphics[width=0.225cm]{Bilder/Sigma.png} \nodepart{lower} \includegraphics[width=0.225cm]{Bilder/sigma.png}};
            }{}%
%-------------- RBF-verdeckt -------------%
            \ifthenelse{\equal{#1}{rbf-h}}{%
                \node[neuron,circle split,inner sep=1pt,font={\fontsize{6pt}{6pt}\selectfont}] (Or) at (0,0)
                    { $\mathbf{||c\mathord{,}x||}$ \nodepart{lower} };
                \node[font={\fontsize{6pt}{6pt}\selectfont}] at ($(Or.lower)-(0,\siz/6)$){\textbf{Gauß}};
            }{}%
%--------------- RBF-Ausgabe -------------%
            \ifthenelse{\equal{#1}{rbf-o}}{%
                \node[neuron,circle split,inner sep=1.3pt] (Or) at (0,0)
                    {\includegraphics[width=0.225cm]{Bilder/Sigma.png} \nodepart{lower} };
                \pic at ($(Or.lower)-(0,\siz/6)$) {ident8};
            }{}%
    }
}


%\usepackage{adjustbox}

\setlength{\parindent}{0pt}             % Entfernt den Einzug der ersten Zeile

\pagestyle{fancy}                       % Ermöglicht benutzerdefinierte Kopf- und Fußzeilen
\lhead{}
\chead{}
\rhead{\thepage}                        % Seitenangabe oben rechts
\cfoot{} 
\renewcommand{\headrulewidth}{0pt}      % Breite des Dekostriches in der Kopfzeile
\renewcommand{\footrulewidth}{0pt}      % Breite des Dekostriches in der Fußzeile

\newcommand{\titel}{Modellierung und Prognose von Strompreisen durch Neuronale Netze}   % Erstellt eine Variable mit dem Titel der Arbeit
\newcommand{\abgabedatum}{!!13. April 2018!!}   % Erstellt eine Variable mit dem Abgabedatum der Arbeit
\newcommand{\name}{Dennis Albrecht}     % Erstellt eine Variable mit dem Namen des Verfassers

\newcommand\blankpage{%                 % Ermöglicht die Erstellung leerer Seiten mit \blankpage
    \null
    \thispagestyle{empty}%
    \newpage%
}

\renewcommand{\cftsecleader}{\textbf{\cftdotfill{\cftdotsep}}} % Füllt den Leerraum der Sectionenangaben im Inhaltsverzeichnis mit Punkten bis zur Seitenzahl (standardmäßig bei article-class nur ab subsections)

%-------------------------------------------
%% Einstellungen für die richtige Nummerierung von Fußnoten und Hyperlinks
%% Quelle: https://tex.stackexchange.com/questions/41626/wrapfigure-footnote-hyperlinks
\makeatletter
\newcommand{\wrapfigfoot}{%
\addtocounter{footnote}{+1}%
\addtocounter{Hfootnote}{+1}%
\global\let\Hy@saved@currentHref\@currentHref%
\hyper@makecurrent{Hfootnote}%
\global\let\Hy@footnote@currentHref\@currentHref%
\global\let\@currentHref\Hy@saved@currentHref%
}
\makeatother
%-------------------------------------------